\documentclass{amsart}
\usepackage[UTF8,fontset=windows]{ctex}
\usepackage[a4paper,margin=24mm]{geometry}
% For dealing with references we use the comment environment
\usepackage{verbatim}
\newenvironment{reference}{\comment}{\endcomment}
%\newenvironment{reference}{}{}
\newenvironment{slogan}{\comment}{\endcomment}
\newenvironment{history}{\comment}{\endcomment}

% For commutative diagrams we use Xy-pic
\usepackage[all]{xy}

% We use 2cell for 2-commutative diagrams.
\xyoption{2cell}
\UseAllTwocells

% We use multicol for the list of chapters between chapters
\usepackage{multicol}

% This is generally recommended for better output
\usepackage{lmodern}
\usepackage[T1]{fontenc}

% For cross-file-references

% Package for hypertext links:
\usepackage{hyperref}

% For any local file, say "hello.tex" you want to link to please

% Theorem environments.
%
\theoremstyle{plain}
\newtheorem{theorem}[subsection]{Theorem}
\newtheorem{proposition}[subsection]{Proposition}
\newtheorem{lemma}[subsection]{Lemma}

\theoremstyle{definition}
\newtheorem{definition}[subsection]{Definition}
\newtheorem{example}[subsection]{Example}
\newtheorem{exercise}[subsection]{Exercise}
\newtheorem{situation}[subsection]{Situation}

\theoremstyle{remark}
\newtheorem{remark}[subsection]{Remark}
\newtheorem{remarks}[subsection]{Remarks}

\numberwithin{equation}{subsection}

% Macros
%
\def\lim{\mathop{\mathrm{lim}}\nolimits}
\def\colim{\mathop{\mathrm{colim}}\nolimits}
\def\Spec{\mathop{\mathrm{Spec}}}
\def\Hom{\mathop{\mathrm{Hom}}\nolimits}
\def\Ext{\mathop{\mathrm{Ext}}\nolimits}
\def\SheafHom{\mathop{\mathcal{H}\!\mathit{om}}\nolimits}
\def\SheafExt{\mathop{\mathcal{E}\!\mathit{xt}}\nolimits}
\def\Sch{\mathit{Sch}}
\def\Mor{\mathop{\mathrm{Mor}}\nolimits}
\def\Ob{\mathop{\mathrm{Ob}}\nolimits}
\def\Sh{\mathop{\mathit{Sh}}\nolimits}
\def\NL{\mathop{N\!L}\nolimits}
\def\CH{\mathop{\mathrm{CH}}\nolimits}
\def\proetale{{pro\text{-}\acute{e}tale}}
\def\etale{{\acute{e}tale}}
\def\QCoh{\mathit{QCoh}}
\def\Ker{\mathop{\mathrm{Ker}}}
\def\Im{\mathop{\mathrm{Im}}}
\def\Coker{\mathop{\mathrm{Coker}}}
\def\Coim{\mathop{\mathrm{Coim}}}

% Boxtimes
%
\DeclareMathSymbol{\boxtimes}{\mathbin}{AMSa}{"02}

%
% Macros for moduli stacks/spaces
%
\def\QCohstack{\mathcal{QC}\!\mathit{oh}}
\def\Cohstack{\mathcal{C}\!\mathit{oh}}
\def\Spacesstack{\mathcal{S}\!\mathit{paces}}
\def\Quotfunctor{\mathrm{Quot}}
\def\Hilbfunctor{\mathrm{Hilb}}
\def\Curvesstack{\mathcal{C}\!\mathit{urves}}
\def\Polarizedstack{\mathcal{P}\!\mathit{olarized}}
\def\Complexesstack{\mathcal{C}\!\mathit{omplexes}}
% \Pic is the operator that assigns to X its picard group, usage \Pic(X)
% \Picardstack_{X/B} denotes the Picard stack of X over B
% \Picardfunctor_{X/B} denotes the Picard functor of X over B
\def\Pic{\mathop{\mathrm{Pic}}\nolimits}
\def\Picardstack{\mathcal{P}\!\mathit{ic}}
\def\Picardfunctor{\mathrm{Pic}}
\def\Deformationcategory{\mathcal{D}\!\mathit{ef}}

\hypersetup{hidelinks}
\widowpenalty=10000\clubpenalty=10000\displaywidowpenalty=10000
\begin{document}
\title[Stacks Project, Tag 089J]{代数空间的弱 Chow 覆盖}
\maketitle
\noindent\textbf{Counting scope.} The primary problem is Tag 089J, marked below. All preceding numbered statements are supporting author context and are not counted as additional problems. This is a local pilot with a temporary ID, not production admission.


\noindent Copyright (C) 2005--2025 Johan de Jong. Stacks Project authors. Source: \href{https://stacks.math.columbia.edu/tag/089J}{Tag 089J}.

\noindent Difficulty H3: 从有限型分离代数空间构造由 H-拟射影概形给出的真满覆盖，不能仅引用概形的 Chow 引理。证明以去掉对角的多重纤维积构造闭包，利用估值提升证明真性，并对几何纤维点数归纳覆盖剩余闭集。核心覆盖对象与归纳装配由本证明完成，故保留为 H3 候选，待独立确认。.

\noindent History: excerpt from revision local cached snapshot; exact chapter hashes; revision unknown, spaces-cohomology.tex, lines 3280--3396. Only reference presentation was adapted. Licensed under GNU FDL 1.2 or later; no invariant sections or cover texts. See accompanying COPYING.

\begin{definition}
\label{morphisms-definition-quasi-projective}
\begin{reference}
\href{https://stacks.math.columbia.edu/bibliography}{Bibliography: EGA (II, Definition 5.3.1)} and \href{https://stacks.math.columbia.edu/bibliography}{Bibliography: H (page 103)}
\end{reference}
Let $f : X \to S$ be a morphism of schemes.
\begin{enumerate}
\item We say $f$ is {\it quasi-projective} if $f$ is of finite type
and there exists an $f$-relatively ample invertible $\mathcal{O}_X$-module.
\item We say $f$ is {\it H-quasi-projective} if there exists
a quasi-compact immersion $X \to \mathbf{P}^n_S$ over $S$ for some
$n$.\footnote{This is not exactly the same as the definition in Hartshorne.
Namely, the definition in Hartshorne (8th corrected printing, 1997) is that
$f$ should be the composition of an open immersion followed by a H-projective
morphism (see Definition \ref{morphisms-definition-projective}), which does not imply
$f$ is quasi-compact. See
Lemma \href{https://stacks.math.columbia.edu/tag/01WA}{lemma H quasi projective open H projective (Tag 01WA)} for
the implication in the other direction.}
\item We say $f$ is {\it locally quasi-projective} if there exists
an open covering $S = \bigcup V_j$ such that each $f^{-1}(V_j) \to V_j$
is quasi-projective.
\end{enumerate}
\end{definition}

\section{Projective morphisms}
\label{morphisms-section-projective}

\noindent
We will use the definition of a projective morphism from \href{https://stacks.math.columbia.edu/bibliography}{Bibliography: EGA}.
The version of the definition with the ``H'' is the one
from \href{https://stacks.math.columbia.edu/bibliography}{Bibliography: H}. The resulting definitions are different. Both are useful.



\begin{lemma}
\label{spaces-morphisms-lemma-refined-valuative-criterion-universally-closed}
Let $S$ be a scheme. Let $f : X \to Y$ and $h : U \to X$ be morphisms of
algebraic spaces over $S$. If
\begin{enumerate}
\item $f$ and $h$ are quasi-compact,
\item $|h|(|U|)$ is dense in $|X|$, and
\end{enumerate}
given any commutative solid diagram
$$
\xymatrix{
\Spec(K) \ar[r] \ar[d] & U \ar[r] & X \ar[d] \\
\Spec(A) \ar[rr] \ar@{-->}[rru] & & Y
}
$$
where $A$ is a valuation ring with field of fractions $K$
\begin{enumerate}
\item[(3)] there exists at most one dotted arrow making the diagram
commute, and
\item[(4)] there exists an extension $K'/K$ of fields, a
valuation ring $A' \subset K'$ dominating $A$ and a morphism
$\Spec(A') \to X$ such that the following diagram commutes
$$
\xymatrix{
\Spec(K') \ar[r] \ar[d] & \Spec(K) \ar[r] & U \ar[r] & X \ar[d] \\
\Spec(A') \ar[r] \ar[rrru] & \Spec(A) \ar[rr] & & Y
}
$$
\end{enumerate}
then $f$ is universally closed. If moreover
\begin{enumerate}
\item[(5)] $f$ is quasi-separated
\end{enumerate}
then $f$ is separated and universally closed.
\end{lemma}

\begin{lemma}
\label{decent-spaces-lemma-bounded-fibres}
Let $S$ be a scheme. Let $X$ be an algebraic space over $S$.
Consider the following conditions on $X$:
\begin{itemize}
\item[] $(\alpha)$ For every $x \in |X|$, the equivalent conditions of
Lemma \href{https://stacks.math.columbia.edu/tag/03JS}{lemma U finite above x (Tag 03JS)}
hold.
\item[] $(\beta)$ For every $x \in |X|$, the equivalent conditions of
Lemma \href{https://stacks.math.columbia.edu/tag/03JU}{lemma R finite above x (Tag 03JU)}
hold.
\item[] $(\gamma)$ For every $x \in |X|$, the equivalent conditions of
Lemma \href{https://stacks.math.columbia.edu/tag/03JV}{lemma UR finite above x (Tag 03JV)}
hold.
\item[] $(\delta)$ The equivalent conditions of
Lemma \href{https://stacks.math.columbia.edu/tag/03JT}{lemma U universally bounded (Tag 03JT)}
hold.
\item[] $(\epsilon)$ The equivalent conditions of
Lemma \href{https://stacks.math.columbia.edu/tag/03IH}{lemma characterize very reasonable (Tag 03IH)}
hold.
\item[] $(\zeta)$ The space $X$ is Zariski locally quasi-separated.
\item[] $(\eta)$ The space $X$ is quasi-separated
\item[] $(\theta)$ The space $X$ is representable, i.e., $X$ is a scheme.
\item[] $(\iota)$ The space $X$ is a quasi-separated scheme.
\end{itemize}
We have
$$
\xymatrix{
& (\theta) \ar@{=>}[rd] & & & &  \\
(\iota) \ar@{=>}[ru] \ar@{=>}[rd] & &
(\zeta) \ar@{=>}[r] &
(\epsilon) \ar@{=>}[r] &
(\delta) \ar@{=>}[r] &
(\gamma) \ar@{<=>}[r] & (\alpha) + (\beta) \\
& (\eta) \ar@{=>}[ru] & & & &
}
$$
\end{lemma}

\begin{lemma}
\label{decent-spaces-lemma-quasi-compact-reasonable-stratification}
Let $S$ be a scheme. Let $W \to X$ be a morphism of a scheme $W$
to an algebraic space $X$ which is flat, locally of finite presentation,
separated, locally quasi-finite with universally bounded fibres. There exist
reduced closed subspaces
$$
\emptyset = Z_{-1} \subset Z_0 \subset Z_1 \subset Z_2 \subset
\ldots \subset Z_n = X
$$
such that with $X_r = Z_r \setminus Z_{r - 1}$ the stratification
$X = \coprod_{r = 0, \ldots, n} X_r$ is characterized by the following
universal property: Given $g : T \to X$ the projection
$W \times_X T \to T$ is finite locally free of degree $r$ if and only if
$g(|T|) \subset |X_r|$.
\end{lemma}

\begin{definition}
\label{morphisms-definition-projective}
Let $f : X \to S$ be a morphism of schemes.
\begin{enumerate}
\item We say $f$ is {\it projective} if $X$ is isomorphic as
an $S$-scheme to a closed subscheme of a projective
bundle $\mathbf{P}(\mathcal{E})$
for some quasi-coherent, finite type $\mathcal{O}_S$-module $\mathcal{E}$.
\item We say $f$ is {\it H-projective} if there exists an integer $n$ and
a closed immersion $X \to \mathbf{P}^n_S$ over $S$.
\item We say $f$ is {\it locally projective} if there exists an open
covering $S = \bigcup U_i$ such that each $f^{-1}(U_i) \to U_i$ is
projective.
\end{enumerate}
\end{definition}

\section*{Primary problem: Tag 089J}
\begin{lemma}
\label{lemma-weak-chow}
Let $A$ be a ring. Let $X$ be an algebraic space over $\Spec(A)$
whose structure morphism $X \to \Spec(A)$ is separated of finite type.
Then there exists a proper surjective morphism $X' \to X$
where $X'$ is a scheme which is H-quasi-projective over $\Spec(A)$.
\end{lemma}

\begin{proof}
Let $W$ be an affine scheme and let $f : W \to X$ be a surjective
\'etale morphism. There exists an integer $d$ such that all geometric
fibres of f have $\leq d$ points (because $X$ is a separated algebraic
hence reasonable, see
Decent Spaces, Lemma \ref{decent-spaces-lemma-bounded-fibres}).
Picking $d$ minimal we get a nonempty open $U \subset X$ such that
$f^{-1}(U) \to U$ is finite \'etale of degree $d$, see
Decent Spaces, Lemma
\ref{decent-spaces-lemma-quasi-compact-reasonable-stratification}.
Let
$$
V \subset W \times_X W \times_X \ldots \times_X W
$$
($d$ factors in the fibre product) be the complement of all the diagonals.
Because $W \to X$ is separated the diagonal $W \to W \times_X W$ is a
closed immersion. Since $W \to X$ is \'etale the diagonal
$W \to W \times_X W$ is an open immersion, see
Morphisms of Spaces, Lemmas
\href{https://stacks.math.columbia.edu/tag/06CR}{lemma etale unramified (Tag 06CR)} and
\href{https://stacks.math.columbia.edu/tag/05W1}{lemma diagonal unramified morphism (Tag 05W1)}.
Hence the diagonals are open and closed subschemes
of the quasi-compact scheme $W \times_X \ldots \times_X W$.
In particular we conclude $V$ is a quasi-compact scheme.
Choose an open immersion $W \subset Y$ with $Y$ H-projective over
$A$ (this is possible as $W$ is affine and of finite type over $A$;
for example we can use
Morphisms, Lemmas
\href{https://stacks.math.columbia.edu/tag/04II}{lemma quasi affine finite type over S (Tag 04II)} and
\href{https://stacks.math.columbia.edu/tag/01WA}{lemma H quasi projective open H projective (Tag 01WA)}).
Let
$$
Z \subset Y \times_A Y \times_A \ldots \times_A Y
$$
be the scheme theoretic image of the composition
$V \to W \times_X \ldots \times_X W \to Y \times_A \ldots \times_A Y$.
Observe that this morphism is quasi-compact since $V$ is quasi-compact
and $Y \times_A \ldots \times_A Y$ is separated.
Note that $V \to Z$ is an open immersion as
$V \to Y \times_A \ldots \times_A Y$ is an immersion, see
Morphisms, Lemma \href{https://stacks.math.columbia.edu/tag/01RG}{lemma quasi compact immersion (Tag 01RG)}.
The projection morphisms give $d$ morphisms $g_i : Z \to Y$.
These morphisms $g_i$ are projective as $Y$ is projective over $A$, see
material in Morphisms, Section \ref{morphisms-section-projective}.
We set
$$
X' = \bigcup g_i^{-1}(W) \subset Z
$$
There is a morphism $X' \to X$ whose restriction to $g_i^{-1}(W)$ is
the composition $g_i^{-1}(W) \to W \to X$.
Namely, these morphisms agree over $V$ hence agree over
$g_i^{-1}(W) \cap g_j^{-1}(W)$ by
Morphisms of Spaces, Lemma \href{https://stacks.math.columbia.edu/tag/084N}{lemma equality of morphisms (Tag 084N)}.
Claim: the morphism $X' \to X$ is proper.

\medskip\noindent
If the claim holds, then the lemma follows by induction on $d$.
Namely, by construction $X'$ is H-quasi-projective over $\Spec(A)$.
The image of $X' \to X$ contains the open $U$ as $V$ surjects onto $U$.
Denote $T$ the reduced induced algebraic space structure on $X \setminus U$.
Then $T \times_X W$ is a closed subscheme of $W$, hence affine.
Moreover, the morphism $T \times_X W \to T$ is \'etale and every geometric
fibre has $< d$ points. By induction hypothesis there exists a proper
surjective morphism $T' \to T$ where $T'$ is a scheme H-quasi-projective
over $\Spec(A)$. Since $T$ is a closed subspace of $X$ we see that
$T' \to X$ is a proper morphism. Thus the lemma follows by taking the
proper surjective morphism $X' \amalg T' \to X$.

\medskip\noindent
Proof of the claim. By construction the morphism $X' \to X$ is separated
and of finite type. We will check conditions (1) -- (4) of
Morphisms of Spaces, Lemma
\ref{spaces-morphisms-lemma-refined-valuative-criterion-universally-closed}
for the morphisms $V \to X'$ and $X' \to X$.
Conditions (1) and (2) we have seen above.
Condition (3) holds as $X' \to X$ is separated (as a morphism whose
source is a separated algebraic space). Thus it suffices to check
liftability to $X'$ for diagrams
$$
\xymatrix{
\Spec(K) \ar[r] \ar[d] & V \ar[d] \\
\Spec(R) \ar[r] & X
}
$$
where $R$ is a valuation ring with fraction field $K$.
Note that the top horizontal map is given by $d$ pairwise distinct
$K$-valued points $w_1, \ldots, w_d$ of $W$. In fact, this
is a complete set of inverse images of the point $x \in X(K)$
coming from the diagram. Since $W \to X$ is surjective,
we can, after possibly replacing $R$ by an extension of valuation rings,
lift the morphism $\Spec(R) \to X$ to a morphism $w : \Spec(R) \to W$, see
Morphisms of Spaces, Lemma
\href{https://stacks.math.columbia.edu/tag/089F}{lemma lift valuation ring through flat morphism (Tag 089F)}.
Since $w_1, \ldots, w_d$ is a complete collection of inverse images of
$x$ we see that $w|_{\Spec(K)}$ is equal to one of them, say $w_i$.
Thus we see that we get a commutative diagram
$$
\xymatrix{
\Spec(K) \ar[r] \ar[d] &  Z \ar[d]_{g_i}\\
\Spec(R) \ar[r]^w & Y
}
$$
By the valuative criterion of properness for the projective
morphism $g_i$ we can lift $w$ to $z : \Spec(R) \to Z$, see
Morphisms, Lemma \href{https://stacks.math.columbia.edu/tag/01WC}{lemma locally projective proper (Tag 01WC)}
and
Schemes, Proposition \href{https://stacks.math.columbia.edu/tag/01KF}{proposition characterize universally closed (Tag 01KF)}.
The image of $z$ is in $g_i^{-1}(W) \subset X'$ and the proof is complete.
\end{proof}
\end{document}
